\subsection*{Problema 41}

\textbf{Enunciado:} \\
Calcular la integral indefinida:
\begin{align*}
I = \int -\dfrac{\sqrt[4]{x}}{\sqrt{x}(1 + \sqrt[4]{x})} \, dx
\end{align*}

\textbf{Solución:} \\
Para resolver la integral, aplicamos la sustitución 
sugerida para eliminar los radicales. 
Sea $u = \sqrt[4]{x}$. De aquí obtenemos que:
\begin{align*}
x &= u^4 \\
dx &= 4u^3 \, du
\end{align*}

Además, expresamos $\sqrt{x}$ en términos de $u$:
\begin{align*}
\sqrt{x} = (x^{1/4})^2 = u^2
\end{align*}

Sustituimos $x$, $\sqrt{x}$ y $dx$ en la integral original:
\begin{align*}
I &= \int -\dfrac{u}{u^2 (1 + u)} (4u^3 \, du)
\end{align*}

Simplificamos la expresión algebraica dentro 
del integrando paso a paso:
\begin{align*}
I &= \int -\dfrac{4u^4}{u^2 (1 + u)} \, du \\
I &= -4 \int \dfrac{u^2}{1 + u} \, du
\end{align*}

Para integrar la función racional impropia, 
realizamos la división polinómica o sumamos 
y restamos $1$ en el numerador:
\begin{align*}
\dfrac{u^2}{1 + u} &= \dfrac{u^2 - 1 + 1}{u + 1} \\
&= u - 1 + \dfrac{1}{1 + u}
\end{align*}

Sustituimos este resultado en la integral:
\begin{align*}
I &= -4 \int \left( u - 1 + \dfrac{1}{1 + u} \right) du
\end{align*}

integramos término a término:
\begin{align*}
I &= -4 \left( \dfrac{u^2}{2} - u + \ln|1 + u| \right) + C \\
I &= -2u^2 + 4u - 4\ln|1 + u| + C
\end{align*}

Finalmente, regresamos a la variable original 
reemplazando $u = \sqrt[4]{x}$ y recordando 
que $u^2 = \sqrt{x}$:
\begin{align*}
I = &-2\sqrt{x} + 4\sqrt[4]{x} \\
&- 4\ln|1 + \sqrt[4]{x}| + C
\end{align*}

\textbf{Resultado final:}
\begin{align*}
\int -\dfrac{\sqrt[4]{x}}{\sqrt{x}(1 + \sqrt[4]{x})} \, dx = \\
-2\sqrt{x} + 4\sqrt[4]{x} - 4\ln(1 + \sqrt[4]{x}) + C
\end{align*}

\begin{figure}[H]
	\centering
	\includegraphics[width=\columnwidth]{semana_05/graficas/SEM05_P41_grafica.png}
	\caption{Representación gráfica de $f(x)$ y su antiderivada $F(x)$ evaluada para la constante de integración $C = 0$ (Problema 41).}
	\label{fig:problema41}
\end{figure}
